Attention costs scale quadratically with sequence length, motivating a long line of approximate and sparse attention shortcuts. Existing locality-sensitive hashing (LSH) approaches index keys by their query embedding at lookup time; none cache a token's \emph{historical} attention pattern keyed by the token's identity for reuse across occurrences. We introduce \ours (\hma), a training-free primitive that builds a per-(layer, head, token-id) lookup table during prefill: when an attention weight crosses a threshold, the key position is stored in the table for that token. At decode, queries for the same token (or for the same n-gram context) compute attention only against stored candidates, attention sinks, and a small local window. We further use the same lookup's hit count to drive \emph{adaptive} \kvcache quantization, allocating more bits to frequently-attended positions. On a planted-structure synthetic stress test, \hma achieves $98\times$ lower attention output error than random selection at matched budget. On \gpttwo-small with \wikitext, \hma matches a sliding-window baseline at $22\%$ smaller effective \kvcache budget (perplexity 90.6 at effective budget 24.8 versus 95.7 at budget 32), and beats H2O and random heavy-hitter baselines by $5{-}25\times$ in perplexity. Hit-count-driven adaptive quantization at 3.9 average bits achieves perplexity 26.4 versus 27.3 for uniform 4-bit, a Pareto improvement at lower memory. \hma also avoids the Top-K aggregation collapse of prior LSH-Top-K methods on a common-word counting probe.
